\notechapter{N-20221030}{Integer Arithmetic: Well-Ordering, Euclidean Algorithm, and the Euler Function}

\section{The well-ordering principle}

The note begins with a concise list of basic facts about integers.  The first is the well-ordering principle:
\begin{theorem}[Well-ordering of \(\mathbb Z_+\)]
Every nonempty subset \(A\subseteq\mathbb Z_+\) has a least element.
\end{theorem}
This statement looks elementary, but it is the engine behind many proofs in number theory.  It is the reason the Euclidean algorithm terminates, the reason minimal-counterexample arguments work, and one of the standard formulations equivalent to induction.

\section{Divisibility, gcd, and lcm}

For integers \(a,b\), with \(a\neq0\), we say
\[
  a\mid b
\]
if there exists \(c\in\mathbb Z\) such that
\[
  b=ac.
\]
The greatest common divisor of nonzero integers \(a,b\) is the positive integer \(d\) such that
\[
  d\mid a,\qquad d\mid b,
\]
and if \(e\mid a\) and \(e\mid b\), then \(e\mid d\).  We write
\[
  d=(a,b).
\]
If \((a,b)=1\), then \(a\) and \(b\) are relatively prime.

Similarly, the least common multiple \(l\) is the positive integer such that
\[
  a\mid l,\qquad b\mid l,
\]
and every common multiple of \(a,b\) is a multiple of \(l\).  For positive \(a,b\),
\[
  (a,b)[a,b]=ab.
\]
The relation is best understood prime-by-prime: gcd takes the smaller exponent, lcm takes the larger exponent.

\section{Division algorithm}

The division algorithm says that for \(a,b\in\mathbb Z\), \(b\neq0\), there exist unique integers \(q,r\) such that
\[
  a=qb+r,\qquad 0\le r<|b|.
\]
Here \(q\) is the quotient and \(r\) is the remainder.  This is not merely long division.  It is the formal foundation of modular arithmetic and the Euclidean algorithm.

\section{Euclidean algorithm}

Starting with \(a,b\neq0\), repeatedly divide:
\[
\begin{aligned}
  a&=q_0b+r_0,\\
  b&=q_1r_0+r_1,\\
  r_0&=q_2r_1+r_2,\\
  &\ \vdots\\
  r_{n-2}&=q_nr_{n-1}+r_n,\\
  r_{n-1}&=q_{n+1}r_n.
\end{aligned}
\]
The last nonzero remainder \(r_n\) is the gcd of \(a,b\).

The reason is the invariant
\[
  (a,b)=(b,r_0)=(r_0,r_1)=\cdots=(r_{n-1},r_n)=r_n.
\]
At each step, replacing \((a,b)\) by \((b,a-qb)\) does not change the set of common divisors.  This is the real idea; the numerical algorithm is just repeated use of this invariant.

\section{Example}

The note gives a Euclidean-algorithm example such as
\[
  57970\quad\text{and}\quad 10353.
\]
The computation proceeds by repeated division until the remainder becomes zero.  The last nonzero remainder is the gcd.  The same backward substitutions produce integers \(x,y\) such that
\[
  ax+by=(a,b).
\]
This is Bézout's identity.

\section{Fundamental theorem of arithmetic}

\begin{theorem}[Fundamental theorem of arithmetic]
Every integer \(n>1\) can be written as a product of primes, and this factorization is unique up to the order of the factors.
\end{theorem}


Every integer \(n>1\) factors uniquely into primes:
\[
  n=p_1^{\alpha_1}p_2^{\alpha_2}\cdots p_k^{\alpha_k}.
\]
The note presents this as an essential property of integers.  Uniqueness is what makes gcd, lcm, and the Euler function computable by exponents.

For example, if
\[
  a=\prod p^{\alpha_p},\qquad b=\prod p^{\beta_p},
\]
then
\[
  (a,b)=\prod p^{\min(\alpha_p,\beta_p)},\qquad
  [a,b]=\prod p^{\max(\alpha_p,\beta_p)}.
\]

\section{The Euler phi function}

The Euler function \(\varphi(n)\) counts the positive integers not exceeding \(n\) that are coprime to \(n\).  If
\[
  n=p_1^{\alpha_1}\cdots p_k^{\alpha_k},
\]
then
\[
  \varphi(n)=n\left(1-\frac1{p_1}\right)\cdots\left(1-\frac1{p_k}\right).
\]
This follows by inclusion--exclusion: remove multiples of each prime divisor of \(n\).

For a prime power,
\[
  \varphi(p^a)=p^a-p^{a-1}.
\]
For a prime \(p\),
\[
  \varphi(p)=p-1.
\]

\section{The broader lesson}

This note is a compact bridge from elementary arithmetic to algebra.  The Euclidean algorithm is not just a computation; it is an algorithmic proof that \(\mathbb Z\) is a Euclidean domain.  Bézout identity says ideals in \(\mathbb Z\) are principal.  Unique factorization says primes control divisibility.  The Euler function prepares the way for the multiplicative group \((\mathbb Z/n\mathbb Z)^\times\), where number theory becomes group theory.

\section{Well-ordering and Euclidean domains}

The well-ordering principle makes the division algorithm possible in \(\mathbb Z\).  The Euclidean algorithm works because each step strictly decreases a nonnegative remainder.  This pattern generalizes to Euclidean domains.

\section{Bezout and principal ideals}

The gcd of \(a\) and \(b\) is the positive generator of the ideal
\[
  (a)+(b)=\{ax+by:x,y\in\mathbb Z\}.
\]
Bezout's identity says the gcd is an integer linear combination of \(a\) and \(b\).

\section{Euler function and unit groups}

Euler's function counts units modulo \(n\):
\[
  \varphi(n)=|(\mathbb Z/n\mathbb Z)^\times|.
\]
Euler's theorem is Lagrange's theorem in this finite group:
\[
  a^{\varphi(n)}\equiv1\pmod n
\]
whenever \(\gcd(a,n)=1\).

\section{Euclidean arithmetic as ideals and unit groups}

Integer arithmetic is organized by order, ideals, and groups.  Well-ordering drives division, Euclidean algorithms compute ideal generators, unique factorization gives valuations, and Euler's theorem is group theory modulo \(n\).

\section{Euclidean domains}

The Euclidean algorithm works in $\mathbb Z$ because there is a division process with decreasing remainder.  A Euclidean domain is an integral domain $R$ with a size function $N$ such that for $a,b\in R$, $b\ne0$, one can write
\[
  a=qb+r,\qquad r=0\text{ or }N(r)<N(b).
\]
The integers and polynomial rings $k[x]$ over a field are Euclidean domains.

In every Euclidean domain, the Euclidean algorithm produces greatest common divisors and Bezout identities.  Thus the elementary gcd algorithm is one example of a general algebraic mechanism.

\section{PID and Bezout domains}

A principal ideal domain is an integral domain in which every ideal is generated by one element.  In a PID, gcds exist because
\[
  (a)+(b)=(d)
\]
for some $d$, and that $d$ is a greatest common divisor of $a$ and $b$.

Bezout's identity
\[
  ax+by=d
\]
means precisely that $d\in(a)+(b)$ and generates the sum ideal.  This reframes gcd as an ideal-theoretic object.

\section{Euler phi as unit counting}

Euler's function counts units in a quotient ring:
\[
  \varphi(n)=|(\mathbb Z/n\mathbb Z)^\times|.
\]
If
\[
  n=\prod_i p_i^{e_i},
\]
then the Chinese remainder theorem gives
\[
  (\mathbb Z/n\mathbb Z)^\times\cong \prod_i(\mathbb Z/p_i^{e_i}\mathbb Z)^\times.
\]
Therefore
\[
  \varphi(n)=n\prod_{p\mid n}\left(1-\frac1p\right).
\]
The familiar formula is a group-size computation.

\section{Multiplicative functions}

A function $f(n)$ is multiplicative if $f(mn)=f(m)f(n)$ whenever $\gcd(m,n)=1$.  Euler's phi, the divisor-counting function $\tau(n)$, and the sum-of-divisors function $\sigma(n)$ are multiplicative.

The reason is again Chinese remaindering or prime factorization: arithmetic data often decomposes independently over primes.  This is the elementary beginning of Euler products and analytic number theory.

\section{Euclidean algorithm and ideals}

The Euclidean algorithm does more than compute a gcd.  It constructs coefficients $x,y$ such that
\[
  ax+by=\gcd(a,b).
\]
The back-substitution process is a constructive proof that the ideal $(a,b)$ is principal.

This construction also gives modular inverses.  If $\gcd(a,n)=1$, then
\[
  ax+ny=1,
\]
so $x$ is the inverse of $a$ modulo $n$.  The original algorithmic procedure is therefore the computational heart of quotient-ring arithmetic.

\section{Euler's theorem as group theory}

Euler's theorem
\[
  a^{\varphi(n)}\equiv1\pmod n
\]
for $\gcd(a,n)=1$ follows from Lagrange's theorem in the finite group $(\mathbb Z/n\mathbb Z)^\times$.  The elementary proof by reduced residue systems is already a group action proof: multiplication by $a$ permutes the reduced residues.

This viewpoint clarifies why the exponent $\varphi(n)$ appears: it is the order of the group of units.
